% ═════════════════════════════════════════════════════════════
% SECTION 6 — Quantifying the Trade-off
% New section. Three subsections: efficiency, relative-improvement
% paradox, marginal sensitivity. Placed after Section 5 (mechanism),
% before Section 7 (cross-model).
%
% Tables used:
%   tab_efficiency_avg          (inline, 6.1)
%   tab_efficiency_ml/ab        (appendix — corrected captions provided)
%   tab_relative_improvement    (inline, 6.2 — corrected caption provided)
%   tab_marginal_sensitivity    (inline, 6.3 — regenerated, all 4 steps)
% Figures: fig_efficiency, fig_relative_improvement, fig_marginal_sensitivity
%
% Refs: sec:mitigation (§4), sec:zone_migration (§5),
%       sec:nash_selection (§3), sec:cross_model (§7),
%       subsec:structural_ceiling (§4.2)
% ═════════════════════════════════════════════════════════════

\section{Quantifying the Trade-off}
\label{sec:tradeoff}

Sections~\ref{sec:mitigation} and~\ref{sec:zone_migration} established
that the achievable long-tail exposure falls with bias and that the
remaining exposure is sourced from progressively deeper, lower-relevance
regions of the pool. This section turns those qualitative findings into
three quantitative, deployment-relevant measures: the efficiency of each
method as an exchange rate between accuracy and fairness
(Section~\ref{subsec:efficiency}); the relative improvement each method
reports over the baseline, which we show to be systematically misleading
at high bias (Section~\ref{subsec:relative}); and the marginal cost of
each successive increment of bias (Section~\ref{subsec:marginal}).

\subsection{Efficiency: the Accuracy--Fairness Exchange Rate}
\label{subsec:efficiency}

We define mitigation efficiency as the APLT@10 gain obtained per
$1\%$ of NDCG@10 sacrificed relative to the baseline --- the exchange
rate at which a method converts accuracy into long-tail exposure. A
higher value indicates a more favourable trade-off.
Table~\ref{tab:efficiency_avg} reports efficiency averaged across the
three architectures; the per-architecture breakdown is deferred to
Appendix~\ref{app:efficiency} (Tables~\ref{tab:efficiency_ml}
and~\ref{tab:efficiency_ab}). Figure~\ref{fig:efficiency} plots the
averaged values.

\begin{table}[htbp]
\centering
\caption{Mitigation efficiency (APLT gain per 1\% NDCG cost) averaged across BPR, NeuMF, and LightGCN. Higher values indicate a more favourable accuracy--fairness trade-off. Smooth xQuAD and FA*IR are statistically indistinguishable in efficiency at every bias level (differences $\le 0.0001$), indicating a shared accuracy--fairness exchange rate; Binary xQuAD is less efficient at low bias but its efficiency rises with $\alpha$.}
\label{tab:efficiency_avg}
\begin{tabular}{clcc}
\toprule
$\alpha$ & Method & MovieLens-1M & Amazon Books \\
\midrule
1.00 & Binary xQuAD & 0.0074 & 0.0095 \\
\rowcolor{gray!6}  & Smooth xQuAD & 0.0136 & 0.0159 \\
 & FA*IR & 0.0137 & 0.0160 \\
\midrule
1.20 & Binary xQuAD & 0.0077 & 0.0104 \\
\rowcolor{gray!6}  & Smooth xQuAD & 0.0134 & 0.0170 \\
 & FA*IR & 0.0134 & 0.0171 \\
\midrule
1.50 & Binary xQuAD & 0.0084 & 0.0136 \\
\rowcolor{gray!6}  & Smooth xQuAD & 0.0134 & 0.0170 \\
 & FA*IR & 0.0134 & 0.0170 \\
\midrule
1.75 & Binary xQuAD & 0.0091 & 0.0139 \\
\rowcolor{gray!6}  & Smooth xQuAD & 0.0132 & 0.0157 \\
 & FA*IR & 0.0133 & 0.0158 \\
\midrule
2.00 & Binary xQuAD & 0.0104 & 0.0135 \\
\rowcolor{gray!6}  & Smooth xQuAD & 0.0132 & 0.0147 \\
 & FA*IR & 0.0131 & 0.0147 \\
\bottomrule
\end{tabular}
\end{table}

\begin{figure}[htbp]
\centering
\includegraphics[width=\linewidth]{template/chapters/figures/fig_efficiency.pdf}
\caption{Mitigation efficiency (APLT gain per $1\%$ NDCG cost) averaged
         across BPR, NeuMF, and LightGCN, for (a) MovieLens-1M and
         (b) Amazon Books. Shaded bands show $\pm 1$ std across
         architectures. Smooth xQuAD and FA*IR are indistinguishable at
         every bias level; Binary xQuAD is less efficient but its
         efficiency rises with $\alpha$.}
\label{fig:efficiency}
\end{figure}

The central observation is that Smooth xQuAD and FA*IR purchase fairness
at \emph{the same exchange rate}. Their efficiencies are equal to within
$0.0001$ at every bias level on both datasets (MovieLens-1M
$\alpha = 1.0$: $0.0136$ vs.\ $0.0137$; Amazon Books $\alpha = 1.0$:
$0.0159$ vs.\ $0.0160$), and the two curves in
Figure~\ref{fig:efficiency} are visually coincident. The two methods
differ in how \emph{much} fairness they buy --- FA*IR consistently
reaches a higher absolute APLT (Section~\ref{subsec:mitigation_effectiveness})
--- but not in the rate at which they trade accuracy for it. This is a
useful deployment conclusion: the choice between Smooth xQuAD and FA*IR
is a choice of operating point along a shared trade-off curve, not a
choice between more- and less-efficient mechanisms.

Binary xQuAD behaves differently in two respects. It is markedly less
efficient at low bias (MovieLens-1M $\alpha = 1.0$: $0.0074$, roughly
half the efficiency of the other two methods), but its efficiency
\emph{rises} with bias, from $0.0074$ at $\alpha = 1.0$ to $0.0104$ at
$\alpha = 2.0$ on MovieLens-1M. This counter-intuitive trend is
explained by the zone analysis of Section~\ref{sec:zone_migration}:
Binary xQuAD promotes only a small number of items and, unlike the
other methods, scarcely draws on the deep pool, so it incurs little of
the escalating relevance cost that Z-3 sourcing imposes on Smooth xQuAD
and FA*IR. Its modest gains thus come at a proportionally smaller
accuracy cost as bias increases.

\paragraph{A normalisation caveat.}
The per-architecture tables in Appendix~\ref{app:efficiency} show
LightGCN achieving the highest nominal efficiency of the three
architectures at most bias levels. This should not be read as LightGCN
making better use of mitigation. Efficiency is normalised by NDCG cost
\emph{expressed as a percentage of the baseline}, and LightGCN has the
largest baseline NDCG of the three models; an equal absolute accuracy
sacrifice therefore registers as a smaller percentage, inflating the
ratio. We resolve this artifact in Section~\ref{sec:cross_model}, where
absolute APLT gains --- which carry no such normalisation --- tell the
opposite story.

\subsection{The Relative-Improvement Paradox}
\label{subsec:relative}

Fairness gains are frequently reported in the literature as relative
improvements over a baseline. Table~\ref{tab:relative_improvement} shows
why this convention becomes actively misleading in the high-bias regime.

\begin{table}[htbp]
\centering
\caption{Relative APLT improvement (\%) at Nash-selected parameters, averaged across BPR, NeuMF, and LightGCN. Binary xQuAD shows a paradoxical increase in relative gain as $\alpha$ rises: the baseline APLT collapses faster than Binary xQuAD's modest absolute gain, inflating the relative figure. The same inflation affects all three methods: every method's relative gain rises with $\alpha$ even though absolute APLT ceilings fall (Section~\ref{subsec:relative}), making relative improvement an unreliable figure of merit under popularity bias.}
\label{tab:relative_improvement}
\begin{tabular}{clcc}
\toprule
$\alpha$ & Method & MovieLens-1M (\%) & Amazon Books (\%) \\
\midrule
1.00 & Binary xQuAD & 41.3\% & 82.4\% \\
\rowcolor{gray!6}  & Smooth xQuAD & 261.1\% & 533.8\% \\
 & FA*IR & 374.4\% & 736.4\% \\
\midrule
1.20 & Binary xQuAD & 81.5\% & 133.9\% \\
\rowcolor{gray!6}  & Smooth xQuAD & 420.2\% & 718.2\% \\
 & FA*IR & 560.7\% & 947.1\% \\
\midrule
1.50 & Binary xQuAD & 184.4\% & 220.8\% \\
\rowcolor{gray!6}  & Smooth xQuAD & 730.5\% & 980.7\% \\
 & FA*IR & 906.5\% & 1231.5\% \\
\midrule
1.75 & Binary xQuAD & 294.9\% & 327.6\% \\
\rowcolor{gray!6}  & Smooth xQuAD & 979.7\% & 1288.2\% \\
 & FA*IR & 1168.2\% & 1564.1\% \\
\midrule
2.00 & Binary xQuAD & 423.1\% & 352.1\% \\
\rowcolor{gray!6}  & Smooth xQuAD & 1248.7\% & 1328.3\% \\
 & FA*IR & 1445.9\% & 1592.2\% \\
\bottomrule
\end{tabular}
\end{table}

\begin{figure}[htbp]
\centering
\includegraphics[width=\linewidth]{template/chapters/figures/fig_relative_improvement.pdf}
\caption{Relative APLT@10 improvement over the baseline (\%), averaged
         across architectures, for (a) MovieLens-1M and (b) Amazon
         Books. Shaded bands show $\pm 1$ std across architectures. All
         methods' relative gains rise steeply with $\alpha$ even though
         their absolute gains shrink, because the baseline denominator
         collapses faster than the numerator.}
\label{fig:relative_improvement}
\end{figure}

Read naively, the table appears to show mitigation becoming dramatically
\emph{more} effective as bias increases: Binary xQuAD's relative APLT
improvement on MovieLens-1M grows more than tenfold, from $41.3\%$ at
$\alpha = 1.0$ to $423.1\%$ at $\alpha = 2.0$, and every method's
relative gain rises monotonically with $\alpha$ on both datasets. This
is the opposite of the structural ceiling established throughout this
chapter, and it is an artifact.

The resolution is the collapsing denominator. Binary xQuAD's
\emph{absolute} APLT gain barely changes across the bias range --- on
MovieLens-1M it moves only from $0.033$ ($0.191 \to 0.224$) at
$\alpha = 1.0$ to $0.049$ ($0.043 \to 0.092$) at $\alpha = 2.0$ --- but
the baseline it is measured against falls from $0.191$ to $0.043$.
Dividing a near-constant numerator by a collapsing denominator
manufactures a tenfold ``improvement'' from what is in fact a roughly
constant, and very small, absolute gain. The same mechanism inflates the
relative figures for Smooth xQuAD and FA*IR, whose relative gains exceed
$1{,}400\%$ at $\alpha = 2.0$ despite their absolute APLT ceilings having
roughly halved.

This is the same denominator effect identified for the Nash score in
Section~\ref{sec:nash_selection}, here in its most extreme form, and it
carries a concrete methodological warning. \emph{Relative} long-tail
improvement is not a safe figure of merit under popularity bias: a
method can report ever-larger percentage gains while its absolute
contribution stagnates or declines, and two methods with very different
absolute effectiveness can report similar relative gains. Evaluations of
fairness mitigation under biased base models should therefore report
absolute APLT alongside, or in preference to, relative improvement.
Throughout this thesis we rely on absolute APLT for exactly this reason.

\subsection{Marginal Sensitivity to Bias}
\label{subsec:marginal}

The final measure asks how the ceiling collapse is distributed across
the bias range: is each increment of $\alpha$ equally costly, or are
some steps more damaging than others? Table~\ref{tab:marginal_sensitivity}
and Figure~\ref{fig:marginal_sensitivity} report the marginal APLT
decline at each of the four bias steps, averaged across architectures.

\begin{table}[htbp]
\centering
\caption{Marginal APLT decline per $\alpha$ step at Nash-selected
parameters, averaged across BPR, NeuMF, and LightGCN. Negative values
indicate a fall in the APLT ceiling between consecutive bias levels; the
largest decline in each row is bolded. FA*IR loses the most at every
step on both datasets. The method declines peak at the $1.2\to1.5$ step
--- where the candidate-pool long-tail fraction drops most sharply
(Table~\ref{tab:pool_lt}) --- and decelerate thereafter, while the
baseline's decline is concentrated in the earliest steps and flattens as
its APLT approaches zero.}
\label{tab:marginal_sensitivity}
\begin{tabular}{llcccc}
\toprule
Step & Dataset & Baseline & Binary xQuAD & Smooth xQuAD & FA*IR \\
\midrule
$1.0\to1.2$ & MovieLens-1M & $-0.0473$ & $-0.0390$ & $-0.0527$ & $\mathbf{-0.0843}$ \\
\rowcolor{gray!6}  & Amazon Books & $-0.0217$ & $-0.0161$ & $-0.0469$ & $\mathbf{-0.0756}$ \\
\midrule
$1.2\to1.5$ & MovieLens-1M & $-0.0364$ & $-0.0323$ & $-0.0748$ & $\mathbf{-0.1104}$ \\
\rowcolor{gray!6}  & Amazon Books & $-0.0218$ & $-0.0226$ & $-0.0838$ & $\mathbf{-0.1224}$ \\
\midrule
$1.5\to1.75$ & MovieLens-1M & $-0.0169$ & $-0.0203$ & $-0.0653$ & $\mathbf{-0.0883}$ \\
\rowcolor{gray!6}  & Amazon Books & $-0.0111$ & $-0.0158$ & $-0.0623$ & $\mathbf{-0.0856}$ \\
\midrule
$1.75\to2.0$ & MovieLens-1M & $-0.0104$ & $-0.0162$ & $-0.0496$ & $\mathbf{-0.0648}$ \\
\rowcolor{gray!6}  & Amazon Books & $-0.0069$ & $-0.0146$ & $-0.0549$ & $\mathbf{-0.0715}$ \\
\bottomrule
\end{tabular}
\end{table}


\begin{figure}[htbp]
\centering
\includegraphics[width=\linewidth]{template/chapters/figures/fig_marginal_sensitivity.pdf}
\caption{Absolute APLT decline at each bias step, averaged across BPR,
         NeuMF, and LightGCN, for (a) MovieLens-1M and (b) Amazon Books.
         Error bars show $\pm 1$ std across architectures. Larger bars
         indicate a faster ceiling collapse at that step. FA*IR loses
         the most at every step, and the method declines peak at the
         $1.2 \to 1.5$ step.}
\label{fig:marginal_sensitivity}
\end{figure}

Two patterns stand out. First, the per-step loss is strictly ordered by
method aggressiveness at every step: FA*IR loses the most APLT per bias
increment, followed by Smooth xQuAD, then Binary xQuAD, with the
baseline losing least in absolute terms. On MovieLens-1M, FA*IR's
largest single-step loss is $0.110$ APLT units (at the $1.2 \to 1.5$
step), against Binary xQuAD's $0.032$ at the same step. The methods that
extract the most long-tail exposure are also those that forfeit the most
as bias tightens --- consistent with the zone-migration mechanism, in
which the aggressive methods have the most deep-pool territory to lose.

Second, the method declines are not linear in $\alpha$ but peak at the
$1.2 \to 1.5$ step and decelerate thereafter (FA*IR on MovieLens-1M:
$0.084 \to 0.110 \to 0.088 \to 0.065$; on Amazon Books the peak is
sharper, $0.076 \to 0.122 \to 0.086 \to 0.072$). This non-monotone,
concave profile matches the candidate-pool evidence of
Section~\ref{sec:bias_validation}: the pool long-tail fraction falls
most steeply in absolute terms over the early-to-middle bias range
(MovieLens-1M $13.3\% \to 9.0\%$ across the $1.2 \to 1.5$ step), where
the reachable Z-2 supply is being exhausted; once that supply is largely
gone, each further bias increment removes less that the methods were
still able to use, and the marginal loss shrinks. The baseline, by
contrast, loses most in the earliest steps and flattens as its APLT
approaches zero, reflecting simple saturation rather than pool
exhaustion. The structural ceiling therefore does not descend at a
constant rate: its steepest descent coincides with the depletion of the
reachable pool, precisely the transition that the zone-migration
crossover in Section~\ref{sec:zone_migration} tracks from the
reranker's side.
